% Research proposal template for phd-compass. Keep phd-compass.sty in the same folder.
% Build: latexmk -pdf proposal.tex. Follow the call's length limit.
\documentclass[10pt]{article}
\usepackage{phd-compass}
\hypersetup{pdftitle={Research proposal - <Name>}, pdfauthor={<Name>}}

\begin{document}

\begin{center}
  {\Large\bfseries <Proposal title>\par}
  \vspace{6pt}
  {\small <Name> \enspace\textbullet\enspace Proposal for <group or programme, institution> \enspace\textbullet\enspace <Month Year>\par}
\end{center}

\section{Background and gap}
<What is known, with citations [1], and the specific gap this project addresses.>

\section{Questions and objectives}
\begin{itemize}
  \item <Question 1, with a measurable objective.>
  \item <Question 2, with a measurable objective.>
\end{itemize}

\section{Approach and methods}
<Methods, data or experiments for each objective; why the host group's facilities and expertise make it feasible.>

\section{Work plan}
\begin{tabular}{@{}p{0.14\textwidth}p{0.80\textwidth}@{}}
  Year 1 & <Milestones> \\
  Year 2 & <Milestones> \\
  Year 3 & <Milestones and thesis writing> \\
\end{tabular}

\section{Risks and alternatives}
<The main risk for each objective and the fallback.>

\section{Fit}
<Your preparation for this project, and how it builds on the group's recent work.>

\section*{References}
\begin{enumerate}[label={[\arabic*]}, leftmargin=20pt, itemsep=1pt]
  \item <Authors, title, venue, year, DOI.>
\end{enumerate}

\end{document}
